\documentclass{presentation}

\title{Άλγεβρα Α Λυκείου}
\subtitle{Η συνάρτηση $y=αx+β$}
\author[Λόλας]{Κωνσταντίνος Λόλας}
\date{\today}

\begin{document}

\frame{\titlepage}

\section{Θεωρία}

\begin{frame}{Ευθεία Λοιπόν}
  Πάμε να θυμηθούμε τα όσα σας έπριξε ο καθηγητής σας στο Γυμνάσιο για την ευθεία! Π.χ.
  \begin{itemize}
    \item Χρειάζεσαι 2 σημεία για να την ορίσεις
    \item Γωνία με άξονα
    \item Κλίση
    \item Τύπος
    \item Σημεία τομής με άξονες
    \item Κλπ κλπ
  \end{itemize}
  Και σαν να μην έφταναν όλα αυτά, τα βάλατε και στη φυσική φέτος!
\end{frame}

\begin{frame}{Πιο σοβαρά...}
  \begin{block}{Ορισμός}
    Γωνία μιας ευθείας ορίζεται η γωνία που σχηματίζεται όταν περιστρέψουμε τον άξονα $x'x$ κατά τη "θετική" φορά, μέχρι να συμπέσει με την ευθεία.
  \end{block}
\end{frame}

\end{document}